\chapter{Kenali ``Panggung'' Otakmu: Primitives vs Ambassadors}

\begin{insightbox}
\textit{``Cinta memang buta --- dan ada alasan ilmiahnya. Saat kamu jatuh cinta, otakmu dibanjiri kimia yang membuatmu tidak bisa berpikir jernih. Tapi dengan memahami ini, kamu bisa mengambil kembali kendali.''}
\end{insightbox}

\section{Dua Sistem dalam Otakmu}

Dr. Stan Tatkin membagi fungsi otak menjadi dua sistem utama yang selalu bekerja --- kadang harmonis, kadang konflik.

\begin{center}
\begin{tikzpicture}[
    box/.style={draw=#1, rounded corners=15pt, fill=#1!10, 
                minimum width=6cm, minimum height=4cm, align=center,
                line width=3pt},
    arrow/.style={<->, line width=3pt, primary}
]
    % Primitives Box
    \node[box=accent] (prim) at (0,0) {};
    \node[font=\Large\bfseries, accent] at (0,1.5) {PRIMITIVES};
    \node[font=\large, accent!80!black] at (0,0.8) {Sistem Primitif};
    \node[align=left, text width=5cm, font=\small] at (0,-0.5) {
        \begin{itemize}[leftmargin=*, itemsep=0.1cm]
            \item Orientasi: Kelangsungan hidup
            \item Kecepatan: Kilat (otomatis)
            \item Mode: Refleks \& memori
            \item Fokus: Bahaya \& ancaman
            \item Respon: Fight, Flight, Freeze
        \end{itemize}
    };
    \node[font=\small\color{accent}] at (0,-2.3) {Cepat | Otomatis | Tanpa sadar};
    
    % Ambassadors Box
    \node[box=sagegreen] (amb) at (9,0) {};
    \node[font=\Large\bfseries, sagegreen] at (9,1.5) {AMBASSADORS};
    \node[font=\large, sagegreen!80!black] at (9,0.8) {Sistem Evolusi};
    \node[align=left, text width=5cm, font=\small] at (9,-0.5) {
        \begin{itemize}[leftmargin=*, itemsep=0.1cm]
            \item Orientasi: Sosial \& kompleks
            \item Kecepatan: Lambat (berpikir)
            \item Mode: Analisis \& empati
            \item Fokus: Hubungan \& makna
            \item Respon: Diplomasi \& solusi
        \end{itemize}
    };
    \node[font=\small\color{sagegreen}] at (9,-2.3) {Lambat | Rasional | Sadar};
    
    % Arrows
    \draw[arrow] (prim.east) -- (amb.west) node[midway, above, font=\small] {Saling mempengaruhi};
    \draw[arrow, bend left=30] (prim.north east) to node[above, font=\small] {Primitives bisa} node[below, font=\small] {override Ambassadors} (amb.north west);
    \draw[arrow, bend right=30] (amb.south west) to node[below, font=\small] {Ambassadors bisa} node[above, font=\small] {regulate Primitives} (prim.south east);
\end{tikzpicture}
\end{center}

\section{Deep Dive: Primitives (Sistem Primitif)}

\subsection{Apa itu Primitives?}

\keyterm{Primitives} adalah bagian otak yang paling tua secara evolusi. Mereka ada untuk satu tujuan: \textbf{membuatmu tetap hidup}.

\begin{praktikbox}[Fungsi Primitives]
\begin{itemize}
    \item \textbf{Survival-focused:} Selalu mencari ancaman
    \item \textbf{Automatic:} Bereaksi tanpa pikir panjang
    \item \textbf{Memory-based:} Menggunakan memori lama untuk menilai situasi sekarang
    \item \textbf{Reflex-driven:} Refleks adalah bosnya
    \item \textbf{Danger-detection:} Selalu waspada
\end{itemize}
\end{praktikbox}

\subsection{Komponen Kunci Primitives}

\begin{center}
\renewcommand{\arraystretch}{1.5}
\begin{tabular}{|>{\raggedright}p{3.5cm}|>{\raggedright}p{5cm}|>{\raggedright\arraybackslash}p{6.5cm}|}
\hline
\rowcolor{accent!20}
\textbf{Struktur} & \textbf{Fungsi} & \textbf{Dalam Hubungan} \\
\hline
\textbf{Amygdala} & Alarm bahaya & Mendeteksi ancaman dalam suara, wajah, atau gerakan pasangan \\
\hline
\textbf{Brain Stem} & Fungsi dasar (napas, detak jantung) & Saat stres, mengontrol respons fisiologis seperti jantung berdebar \\
\hline
\textbf{Cerebellum} & Koordinasi motorik & Gerakan tiba-tiba saat marah atau ketakutan \\
\hline
\textbf{Basal Ganglia} & Kebiasaan dan reward & Membuatmu ``kecanduan'' pasangan melalui dopamine \\
\hline
\textbf{Hippocampus} (bagian) & Memori emosional & Mengingat pengalaman traumatis masa lalu \\
\hline
\end{tabular}
\end{center}

\subsection{Bagaimana Primitives Bekerja dalam Dating?}

Saat kamu melihat seseorang yang menarik:

\begin{enumerate}
    \item \textbf{Far Visual System} (bagian dari primitives) mendeteksi gerakan, bentuk, sinyal atraksi
    \item \textbf{Amygdala} memutuskan: ``Aman atau berbahaya?''
    \item Jika ``aman'', \textbf{Basal Ganglia} rilis dopamine --- sensasi ``suka''
    \item Jika ada red flag tapi \textbf{familiar}, primitives mungkin bilang ``Aman!'' karena familiar = aman untuk otak primitif
\end{enumerate}

\begin{warningbox}[Bahaya Primitives dalam Hubungan]
\begin{itemize}
    \item \textbf{False alarm:} Amygdala bisa melihat ``ancaman'' padahal pasanganmu hanya capek
    \item \textbf{Familiar trap:} Kamu tertarik pada orang yang seperti orang tua yang menyakitimu (karena familiar)
    \item \textbf{Overreaction:} Respon fight/flight yang berlebihan saat konflik kecil
\end{itemize}
\end{warningbox}

\section{Deep Dive: Ambassadors (Sistem Evolusi)}

\subsection{Apa itu Ambassadors?}

\keyterm{Ambassadors} adalah bagian otak yang lebih baru secara evolusi. Mereka adalah ``diplomat'' yang bisa:

\begin{praktikbox}[Fungsi Ambassadors]
\begin{itemize}
    \item \textbf{Nuanced judgment:} Menilai situasi secara kompleks, tidak hitam-putih
    \item \textbf{Empathy:} Merasakan apa yang dirasakan orang lain
    \item \textbf{Delayed gratification:} Menunda kesenangan untuk tujuan jangka panjang
    \item \textbf{Planning:} Merencanakan dan memprediksi
    \item \textbf{Social regulation:} Mengontrol ekspresi sosial agar sesuai norma
\end{itemize}
\end{praktikbox}

\subsection{Komponen Kunci Ambassadors}

\begin{center}
\renewcommand{\arraystretch}{1.5}
\begin{tabular}{|>{\raggedright}p{3.5cm}|>{\raggedright}p{5cm}|>{\raggedright\arraybackslash}p{6.5cm}|}
\hline
\rowcolor{sagegreen!20}
\textbf{Struktur} & \textbf{Fungsi} & \textbf{Dalam Hubungan} \\
\hline
\textbf{Prefrontal Cortex} & Pengambilan keputusan, kontrol impuls & Menahan diri tidak menyakiti pasangan saat marah \\
\hline
\textbf{Anterior Cingulate} & Deteksi error, empati & Menyadari ``Aku salah'' dan memahami perasaan pasangan \\
\hline
\textbf{Orbital Frontal Cortex} & Penilaian sosial, kontrol emosi & Membaca ekspresi wajah pasangan dengan akurat \\
\hline
\textbf{Hippocampus} (utama) & Memori eksplisit, konteks & Mengingat detail tentang pasangan (Love Maps) \\
\hline
\textbf{Insula} & Kesadaran tubuh, intuisi & ``Merasakan'' bahwa sesuatu salah dalam hubungan \\
\hline
\end{tabular}
\end{center}

\section{Konflik Antara Primitives dan Ambassadors}

\subsection{Skenario: Lihat Mantan di Media Sosial}

\begin{center}
\begin{tikzpicture}[
    scenario/.style={draw=warmgray, rounded corners=10pt, fill=warmgray!10,
                     minimum width=4cm, minimum height=2.5cm, align=center},
    response/.style={draw=#1, rounded corners=5pt, fill=#1!10,
                     minimum width=4cm, minimum height=1.5cm, align=center}
]
    % Trigger
    \node[scenario] (trigger) at (0,0) {\textbf{Trigger}\\\small Melihat foto\\mantan dengan\\pasangan baru};
    
    % Primitives Response
    \node[response=accent] (prim) at (-5,-4) {\textbf{Primitives}\\\small ``Ancaman!''\\Jantung berdebar\\Merasa sakit hati\\Ingin stalking};
    
    % Ambassadors Response
    \node[response=sagegreen] (amb) at (5,-4) {\textbf{Ambassadors}\\\small ``Ini normal''\\Dia bukan ancaman\\Aku punya hidupku\\Tutup aplikasi};
    
    % Final Action
    \node[response=primary] (action) at (0,-7) {\textbf{Action}\\\small Tergantung mana yang menang};
    
    % Arrows
    \draw[->, line width=2pt, accent] (trigger) -- (prim);
    \draw[->, line width=2pt, sagegreen] (trigger) -- (amb);
    \draw[->, line width=2pt] (prim) -- (action);
    \draw[->, line width=2pt] (amb) -- (action);
\end{tikzpicture}
\end{center}

\subsection{Kapan Ambassadors Offline?}

Ambassadors \textbf{tidak berfungsi} saat:
\begin{enumerate}
    \item \textbf{Flooding:} Jantung berdebar $>$100 bpm, darah naik
    \item \textbf{Sleep deprivation:} Kurang tidur
    \item \textbf{Drunk/High:} Terpengaruh zat
    \item \textbf{Extreme hunger:} Sangat lapar
    \item \textbf{Fog of infatuation:} Dopamine tinggi
\end{enumerate}

\begin{tipbox}[Cara Membangunkan Ambassadors]
Saat kamu merasa ``dikuasai'' primitives:
\begin{enumerate}
    \item \textbf{Grounding:} Rasakan kontak kaki dengan lantai
    \item \textbf{Labeling:} Ucapkan ``Ini hanya amydala-ku yang bereaksi''
    \item \textbf{Breathing:} Napas dalam 4-7-8 (tarik 4, tahan 7, hembus 8)
    \item \textbf{Cooling:} Kompres air dingin di wajah (mengaktifkan dive reflex)
    \item \textbf{Time-out:} Beri diri waktu 20 menit untuk turun adrenalin
\end{enumerate}
\end{tipbox}

\section{Fog of Infatuation: Kabut Kegilaan Cinta}

\subsection{Neurochemical Cocktail}

Saat jatuh cinta, otakmu dipenuhi koktail kimia:

\begin{center}
\begin{tikzpicture}[
    chem/.style={draw=#1, line width=2pt, rounded corners=10pt, 
                 fill=#1!10, minimum width=4cm, minimum height=2.5cm, align=center},
    arrow/.style={->, line width=1.5pt}
]
    % Top row
    \node[chem=secondary] (dop) at (0,0) {\textbf{Dopamine}\\\small Pleasure/Reward\\\faArrowUp \ Euforia};
    \node[chem=accent] (nora) at (5,0) {\textbf{Norepinephrine}\\\small Excitement/Focus\\\faArrowUp \ Jantung berdebar};
    \node[chem=warmgray] (sero) at (10,0) {\textbf{Serotonin}\\\small Mood regulation\\\faArrowDown \ Obsesi};
    
    % Bottom row
    \node[chem=softpink] (test) at (2.5,-3.5) {\textbf{Testosterone}\\\small Libido\\Turun (pria), Naik (wanita)};
    \node[chem=sagegreen] (oxy) at (7.5,-3.5) {\textbf{Oxytocin}\\\small Bonding\\\faArrowUp \ Keterikatan};
    
    % Effects
    \node[draw=primary, rounded corners=10pt, fill=primary!10, 
          minimum width=12cm, minimum height=2cm, align=center] (effect) at (5,-6.5) {
        \textbf{Efek Gabungan:} Kurang tidur, kurang makan, obsesi berpikir tentang dia,\\
        mengabaikan red flags, merasa ``ini yang terbaik!''
    };
    
    % Arrows
    \draw[arrow, secondary] (dop) -- (test);
    \draw[arrow, accent] (nora) -- (oxy);
    \draw[arrow, warmgray] (sero) -- (oxy);
    \draw[arrow, softpink] (test) -- (effect);
    \draw[arrow, sagegreen] (oxy) -- (effect);
\end{tikzpicture}
\end{center}

\subsection{Fakta Mengejutkan tentang Serotonin}

\begin{warningbox}[The Serotonin Drop]
\begin{itemize}
    \item Saat jatuh cinta, \textbf{serotonin turun 40\%}
    \item Level yang sama dengan penderita \textbf{OCD} (Obsessive Compulsive Disorder)
    \item Ini sebabnya kamu \textbf{tidak bisa berhenti memikirkan} dia
    \item Ini juga sebabnya kamu \textbf{obsesif} --- cek hp terus, stalking media sosial
\end{itemize}

\textit{``Crazy in love'' bukan hanya metafora --- ini kondisi neurokimia yang nyata!}
\end{warningbox}

\subsection{Durasi Fog of Infatuation}

Kabut ini \textbf{tidak abadi}:
\begin{itemize}
    \item \textbf{3-6 bulan:} Puncak euforia
    \item \textbf{6-12 bulan:} Mulai turun
    \item \textbf{12-24 bulan:} Kembali ke normal
    \item \textbf{Setelah itu:} Attachment (oxytocin/vasopressin) mengambil alih
\end{itemize}

\begin{insightbox}
\textbf{Implikasi:} Jangan ambil keputusan besar (menikah, pindah bareng, investasi besar) dalam 12-18 bulan pertama hubungan. Tunggu sampai fog hilang untuk melihat realitas dengan jelas.
\end{insightbox}

\section{Visual System: Far vs Near}

Otak punya dua jalur visual yang berbeda:

\begin{center}
\renewcommand{\arraystretch}{1.6}
\begin{tabular}{|>{\raggedright}p{3cm}|>{\raggedright}p{5.5cm}|>{\raggedright\arraybackslash}p{5.5cm}|}
\hline
\rowcolor{primary!20}
\textbf{Aspek} & \textbf{Far Visual System} & \textbf{Near Visual System} \\
\hline
\textbf{Alias} & Dorsal stream (``where/how'') & Ventral stream (``what'') \\
\hline
\textbf{Fungsi utama} & Deteksi gerakan, lokasi, arah & Detail wajah, ekspresi, identifikasi \\
\hline
\textbf{Resolusi} & Rendah (blur) & Tinggi (tajam) \\
\hline
\textbf{Hubungan dengan} & \textbf{Primitives} & \textbf{Ambassadors} \\
\hline
\textbf{Jarak optimal} & $>$ 1 meter & $<$ 1 meter \\
\hline
\textbf{Deteksi} & Bahaya, atraksi kasar & Emosi, niat, kebohongan \\
\hline
\textbf{Dalam dating} & ``Orang itu hot!'' (dari kejauhan) & ``Dia bohong'' (dari dekat) \\
\hline
\end{tabular}
\end{center}

\subsection{Mengapa ``Cinta Pada Pandangan Pertama'' Menipu}

Far Visual System:
\begin{itemize}
    \item Melihat gerakan pinggul, cara berjalan, siluet
    \item \textbf{Mengaktifkan Primitives}
    \item Memberi sinyal ``Ini orang yang menarik''
    \item Tapi \textbf{tidak bisa menilai} karakter, kejujuran, kompatibilitas
\end{itemize}

\begin{tipbox}[Gunakan Near Visual System dengan Benar]
Untuk benar-benar mengenal seseorang:
\begin{itemize}
    \item Duduk \textbf{berhadapan} dalam jarak $<$ 1 meter
    \item Pastikan \textbf{PENCAYAAN CUKUP} (jangan tempat gelap)
    \item Lihat ke \textbf{mata} --- pupil, micro-expression
    \item Perhatikan \textbf{asimetri wajah} --- alkohol/membuat wajah terlihat lebih simetris/simetris = menarik
    \item Lakukan ini \textbf{berulang kali} dalam berbagai situasi
\end{itemize}
\end{tipbox}

\section{Sistem Saraf dalam Relasi: Polyvagal Theory}

Stephen Porges mengembangkan teori tentang tiga mode respons saraf kita dalam interaksi sosial.

\begin{center}
\begin{tikzpicture}[
    mode/.style={draw=#1, line width=3pt, rounded corners=15pt,
                 fill=#1!10, minimum width=13cm, minimum height=2cm,
                 align=left, text width=12cm}
]
    \node[mode=sagegreen] (safe) at (0,4) {
        \textbf{\large \textcolor{sagegreen}{SAFETY MODE (Mode Aman / Ventral Vagal)}}\\[0.2cm]
        \begin{itemize}[leftmargin=*, itemsep=0.05cm]
            \item Kontrol napas normal, kontak mata mudah, suara modulasi
            \item Sosial engagement aktif: bisa mendengarkan, berempati, berpikir jernih
            \item \textit{Otak berfungsi optimal}
        \end{itemize}
    };
    
    \node[mode=accent] (danger) at (0,1) {
        \textbf{\large \textcolor{accent}{DANGER MODE (Mode Bahaya / Sympathetic)}}\\[0.2cm]
        \begin{itemize}[leftmargin=*, itemsep=0.05cm]
            \item Fight or Flight: mulut kering, tangan dingin/berkeringat, jantung berdebar
            \item Suara cepat/tinggi, gerakan tajam, blushing
            \item \textit{Fungsi sosial terganggu, sulit berpikir jernih}
        \end{itemize}
    };
    
    \node[mode=warmgray] (threat) at (0,-2) {
        \textbf{\large \textcolor{warmgray}{LIFE THREAT (Ancaman Hidup / Dorsal Vagal)}}\\[0.2cm]
        \begin{itemize}[leftmargin=*, itemsep=0.05cm]
            \item Freeze/Collapse: pusing, mual, lemas, suara sangat pelan
            \item ``Blank out'', dissociation, tidak bisa berbicara
            \item \textit{Shutdown total --- sistem ``pura-pura mati''}
        \end{itemize}
    };
    
    \draw[->, line width=4pt, sagegreen] (-6.5,6) -- (safe.west);
    \draw[->, line width=2pt] (safe.south) -- (danger.north) node[midway, right, font=\small] {Stres berlanjut};
    \draw[->, line width=2pt] (danger.south) -- (threat.north) node[midway, right, font=\small] {Stres ekstrem};
\end{tikzpicture}
\end{center}

\subsection{Co-Regulation: Menenangkan Satu Sama Lain}

\begin{insightbox}
Manusia adalah satu-satunya mamalia yang \textbf{tidak bisa} menenangkan diri sendiri dengan efektif. Kita butuh \keyterm{co-regulation} --- saling menenangkan dengan interaksi sosial.
\end{insightbox}

\textbf{Tanda-tanda butuh co-regulation:}
\begin{itemize}
    \item Jantung berdebar
    \item Napas pendek/cepat
    \item Otot tegang
    \item Pikiran berkecamuk
    \item Ingin menjauh atau menyerang
\end{itemize}

\textbf{Cara co-regulate dengan pasangan:}
\begin{enumerate}
    \item \textbf{Kontak fisik:} Pegang tangan, peluk (jika nyaman)
    \item \textbf{Kontak mata lembut:} Soft eye contact
    \item \textbf{Suara rendah dan lambat:} ``Aku di sini'', ``Kita aman''
    \item \textbf{Breathing together:} Napas bersamaan
    \item \textbf{Validation:} ``Aku mengerti kamu stres''
\end{enumerate}

\section{Familiarity Effect: Jebakan yang Dikenal}

Otak tertarik pada yang \textbf{familiar} --- tapi familiar $\neq$ sehat.

\subsection{Mere Exposure Effect}

Semakin sering kamu melihat sesuatu, semakin kamu menyukainya.

\begin{praktikbox}[Efek Familiar dalam Hubungan]
\textbf{Familiar yang Sehat:}
\begin{itemize}
    \item Orang yang menghargai hubungan dan terbuka
    \item Pola komunikasi jujur dan langsung
    \item Rasa aman tanpa drama
    \item Konflik yang diselesaikan dengan baik
\end{itemize}

\textbf{Familiar yang Tidak Sehat:}
\begin{itemize}
    \item Orang seperti orang tua yang tidak stabil
    \item Drama, chaos, roller-coaster emosi
    \item Sikap dingin/menghindar (jika kamu Island)
    \item Sikap clingy/cemas (jika kamu Wave)
\end{itemize}
\end{praktikbox}

\begin{warningbox}[Trauma Bonding]
Jika kamu dibesarkan dalam lingkungan yang tidak stabil, kamu mungkin merasa ``nyaman'' dengan pasangan yang juga tidak stabil. Ini disebut \textbf{trauma bonding} --- hubungan berbasis trauma yang sama.

\textbf{Tanda-tanda:}
\begin{itemize}
    \item Merasa ``hidup'' hanya saat ada drama
    \item Pasangan yang baik-baik saja terasa membosankan
    \item Selalu tertarik pada orang yang ``perlu diselamatkan''
    \item Pola hubungan berulang yang tidak memuaskan
\end{itemize}
\end{warningbox}

\section{Action Steps: Praktik Mingguan}

\begin{exercisebox}[Program 7 Hari: Mengenali Wiring-mu]

\textbf{Hari 1-2: Mindful Breathing (5 menit)}
\begin{itemize}
    \item Duduk nyaman, punggung lurus
    \item Tarik napas 4 detik $\rightarrow$ Tahan 4 detik $\rightarrow$ Hembusk 6 detik
    \item Fokus hanya pada napas
    \item Jika pikiran melayang, kembalikan tanpa menghakimi
\end{itemize}

\textbf{Hari 3-4: Body Scan saat Dating/Berpasangan}
\begin{itemize}
    \item Sebelum bertemu pasangan, tutup mata
    \item Scan dari kepala ke kaki: di mana ketegangan?
    \item Catat di jurnal: ``Saat akan bertemu dia, aku merasa...''
\end{itemize}

\textbf{Hari 5-6: Deteksi Fog}
\begin{itemize}
    \item Apakah jantungku berdebar kencang memikirkan dia?
    \item Apakah aku sulit tidur/makan karena memikirkannya?
    \item Apakah aku mengabaikan kekurangan yang jelas?
    \item Jika ya, tulis: ``Ini hanya kimia. Aku perlu waktu untuk melihat dengan jelas.''
\end{itemize}

\textbf{Hari 7: Reflection}
\begin{itemize}
    \item Review jurnal minggu ini
    \item Identifikasi: Kapan Primitives-mu bereaksi berlebihan?
    \item Kapan Ambassadors-mu berhasil mengontrol?
    \item Apa strategi yang paling efektif untukmu?
\end{itemize}
\end{exercisebox}

\begin{tipbox}[Quick Reference Card]
\textbf{Saat kamu merasa ``dikuasai'' emosi:}
\begin{enumerate}
    \item \textbf{STOP} --- jangan reaksi dulu
    \item \textbf{BREATHE} --- 3 napas dalam
    \item \textbf{LABEL} --- ``Ini amydala-ku''
    \item \textbf{GROUND} --- Rasakan kaki di lantai
    \item \textbf{WAIT} --- 20 menit sebelum merespons
\end{enumerate}
\end{tipbox}
